%% 该模板修改自《计算机学报》latex 模板
%% 主要是将双栏改成单栏，去掉了部分计算机学报标识；
%% 源文件自：https://www.overleaf.com/latex/templates/latextemplet-cjc-xelatex/ybmmymncrrmw
%% 
%%
%% This is file `CjC_template_tex.tex',
%% is modified by Zhi Wang (zhiwang@ieee.org) based on the template 
%% provided by Chinese Journal of Computers (http://cjc.ict.ac.cn/).
%%
%% This version is capable with Overleaf (XeLaTeX).
%%
%% Update date: 2026/09/10
%% -------------------------------------------------------------------
%% Copyright (C) 2016--2026 
%% -------------------------------------------------------------------
%% This file may be distributed and/or modified under the
%% conditions of the LaTeX Project Public License, either version 1.3c
%% of this license or (at your option) any later version.
%% The latest version of this license is in
%%    https://www.latex-project.org/lppl.txt
%% and version 1.3c or later is part of all distributions of LaTeX
%% version 2008 or later.
%% -------------------------------------------------------------------

\documentclass[10.5pt,compsoc,UTF8]{CjC}
\usepackage{CTEX}
\usepackage{graphicx}
\usepackage{footmisc}
\usepackage{subfigure}
\usepackage{url}
\usepackage{multirow}
\usepackage{multicol}
\usepackage[noadjust]{cite}
\usepackage{amsmath,amsthm}
\usepackage{amssymb,amsfonts}
\usepackage{booktabs}
\usepackage{color}
\usepackage{ccaption}
\usepackage{booktabs}
\usepackage{float}
\usepackage{fancyhdr}
\usepackage{caption}
\usepackage{xcolor,stfloats}
\usepackage{comment}
\setcounter{page}{1}
\graphicspath{{figures/}}
\usepackage{cuted}%flushend,
\usepackage{captionhack}
\usepackage{epstopdf}
\usepackage{gbt7714}

%===============================%

\headevenname{\mbox{\quad} \hfill  \mbox{\zihao{-5}{ \hfill 数字图像处理  } \hspace {50mm} \mbox{2026 年 10 月}}}%
\headoddname{李建辉，苏威贤，王妍，李定一 \hfill 低剂量 CT 图像质量重建与超分辨率增强 }%

%footnote use of *
\renewcommand{\thefootnote}{\fnsymbol{footnote}}
\setcounter{footnote}{0}
\renewcommand\footnotelayout{\zihao{5-}}

\newtheoremstyle{mystyle}{0pt}{0pt}{\normalfont}{1em}{\bf}{}{1em}{}
\theoremstyle{mystyle}
\renewcommand\figurename{figure~}
\renewcommand{\thesubfigure}{(\alph{subfigure})}
\newcommand{\upcite}[1]{\textsuperscript{\cite{#1}}}
\renewcommand{\labelenumi}{(\arabic{enumi})}
\newcommand{\tabincell}[2]{\begin{tabular}{@{}#1@{}}#2\end{tabular}}
\newcommand{\abc}{\color{white}\vrule width 2pt}
\renewcommand{\bibsection}{}
\makeatletter
\renewcommand{\@biblabel}[1]{[#1]\hfill}
\makeatother
\setlength\parindent{2em}
%\renewcommand{\hth}{\begin{CJK*}{UTF8}{gbsn}}
%\renewcommand{\htss}{\begin{CJK*}{UTF8}{gbsn}}


\begin{document}

\hyphenpenalty=50000
\makeatletter
\newcommand\mysmall{\@setfontsize\mysmall{7}{9.5}}
\newenvironment{tablehere}
  {\def\@captype{table}}

\let\temp\footnote
\renewcommand \footnote[1]{\temp{\zihao{-5}#1}}


\thispagestyle{plain}%
\thispagestyle{empty}%
\pagestyle{CjCheadings}

% \begin{table*}[!t]
\vspace {-13mm}


\onecolumn
\zihao{5-}\noindent 李建辉，苏威贤，王妍，李定一 \hfill DIP \& 低剂量 CT 图像质量重建与超分辨率增强 \hfill 2026 年 10 月\\
\noindent\rule[0.25\baselineskip]{\textwidth}{1pt}

% \onecolumn
% \zihao{5-}\noindent 第??卷\quad 第?期 \hfill 计\quad 算\quad 机\quad 学\quad 报\hfill Vol. ??  No. ?\\
% \zihao{5-}
% 20??年?月 \hfill CHINESE JOURNAL OF COMPUTERS \hfill ???. 20??\\ 
% \noindent\rule[0.25\baselineskip]{\textwidth}{1pt}

{
\centering
\vspace {11mm}
{\zihao{2} \heiti 低剂量 CT 图像质量重建与超分辨率增强 }
%题目（中英文题目一致) 字体为 2 号黑体 (全文除特别声明外，外文统一用 Times New Roman)
\vskip 5mm

{\zihao{4}\fangsong 李建辉\quad  苏威贤 \quad 王妍 \quad 李定一}
%字体为 3 号仿宋 * 作者
% \footnote{\noindent \zihao{6} 
% % 收稿日期：\quad \quad -\quad -\quad ；最终修改稿收到日期：\quad \quad -\quad -\quad .*投稿时不填写此项*. 本课题得到… …基金中文完整名称(No.项目号)、… …基金中文完整名称(No.项目号)、… … 基金中文完整名称(No.项目号)资助.
% \textsf{作者名1}，学号，学位(或目前学历),主要研究领域为*****、****. E-mail: **************.\textsf{作者名2}，学号，学位(或目前学历),主要研究领域为*****、****. E-mail: **************. \textsf{作者名3}，学号，学位(或目前学历),主要研究领域为*****、****. E-mail: **************.}

\vspace {5mm}
% \zihao{6}{\textsf{论文定稿后，作者署名、单位无特殊情况不能变更。若变更，须提交签章申请，国家为中国可以不写，省会城市不写省名，其他国家必须写国家名。}}
}

\vskip 5mm
%%%%%%%%%%%%%%%%%%%%%%%%%%摘要%%%%%%%%%%%%%%%%%%%%%%%%%%%%%
\zihao{5}{
\setlength{\baselineskip}{16pt}\selectfont{
\noindent {\heiti 摘\quad 要\quad }
*中文摘要内容置于此处(英文摘要中要有这些内容)，字体为5号宋体。摘要贡献部分，要有数据支持，不要出现``...大大提高''、``...显著改善''等描述，正确的描述是``比{\ldots}提高X{\%}''、``在{\ldots}上改善X{\%}''。
\par}}
\vspace {5mm}

\zihao{5}{\noindent
{\heiti 关键词 \quad }{*关键词（中文关键字与英文关键字对应且一致，应有5-7个关键词)；关键词；关键词*  }
}
% \zihao{5-}{\heiti 中图法分类号\quad } TP\rm{\quad \quad \quad     }
% {\heiti DOI号:\quad } *投稿时不提供DOI号


\vskip 7mm

% \begin{center}
% \zihao{3}{ \heiti Title *（中英文题目一致)字体为4号Times New Roman,加粗* Title}\\
% \vspace {5mm}
% \zihao{5}{ NAME Name-Name$^{1)}$ NAME Name$^{2)}$ NAME Name-Name$^{3)}$ *字体为5号Times new Roman*Name}\\
% \vspace {2mm}
% \zihao{5-}{{$^{1)}$(Department of ****, University, City ZipCode, China) *字体为小5号Times new Roman* Depart.Correspond}}

% \zihao{5-}{{$^{2)}$(Department of ****, University, City ZipCode)*中国不写国家名*}}

% \zihao{5-}{{$^{3)}$(Department of ****, University, City ZipCode, country)*外国写国家名*}}

% \end{center}
%%%%%%%%%%%%%%%%%%%%英文摘要%%%%%%%%%%%%%%%%%%%%%%%%%%%%%
\zihao{5}{
\setlength{\baselineskip}{18pt}\selectfont{
{\bf Abstract}\quad (\textbf{500英文单词，内容包含中文摘要的内容}).
字体为Times new Roman,字号5号* 
\zihao{5}{\noindent Do not modify the amount of space before and after the artworks. One- or two-column format artworks are preferred. and Tables, create a new break line and paste the resized artworks where desired. Do not modify the amount of space before and after the artworks. One- or two-column format artworks are preferred. All Schemes, Equations, Figures, and Tables should be mentioned in the text consecutively and numbered with Arabic numerals, and appear below where they are mentioned for the first time in the main text. To insert Schemes, Equations, Figures, and Tables, create a new break line and paste the resized artworks where desired. Do not modify the amount of space before and after the artworks. One- or two-column format artworks are preferred.Do not modify the amount of space before and after the artworks. One- or two-column format artworks are preferred. and Tables, create a new break line and paste the resized artworks where desired. %Do not modify the amount of space before and after the artworks. One- or two-column format artworks are preferred. All Schemes, Equations, Figures, and Tables should be mentioned in the text consecutively and numbered with Arabic numerals, and appear below where they are mentioned for the first time in the main text.
\par}}
\vspace {5mm}

{{\bf Keywords}\quad 中文关键字与英文关键字对应且一致，\textbf{不要用英文缩写});
key word; key word; key word* *字体为5号Times new Roman * Key words\par}}



%%%%%%%%%%%%%%%%%%%%%%%%%%%%%%%%%%%%%%
\zihao{5}
\vskip 10mm
% \begin{multicols}{1}


%%%%%%%%%%%%%%%%%%%%%%%%%%%%%%%%%%%%%%%%%%%%%%%%%%%%%%%%%%%%%%%%%%%%%%%
%%%%%%%%%%%%%%%%%%%%%%%%%%%%%%%%%%%%%%%%%%%%%%%%%%%%%%%%%%%%%%%%%%%%%%
%%%%%%%%%%%%%%正文%%%%%%%%%%%%%%%%%%%%%%%%%%%%%%%
\section{\heiti 引言}
%介绍低剂量 CT 的噪声问题、去噪与细节保留的矛盾，以及本实验的目标和主要工作。
\section{\heiti 相关工作}
%2.1 传统空间域去噪；2.2 小波域去噪；2.3 图像插值增强。简要比较各类方法的特点，为后续方法选择作铺垫。
\section{\heiti 实验方法与处理流程}
%3.1 总体流程；3.2 DICOM 读取、切片配对、HU 截断与归一化；3.3 均值、高斯、自适应中值、NLM、小波软阈值五种去噪方法；3.4 最近邻、双线性、双三次三种 4 倍插值方法；3.5 PSNR、SSIM 指标与计算条件。
\section{\heiti 实验设置与结果分析}
%4.1 数据选择、运行环境与参数；4.2 五种去噪结果对比；4.3 最优去噪结果的三种插值对比；4.4 “先去噪后插值”与“先插值后去噪”对比；4.5 综合方案选择与局限性分析。
\section{\heiti 结论}
%5 结论概括本次实验中表现较好的方法、处理顺序的影响，以及传统插值无法恢复真实缺失高频细节的限制。结论限定在实际测试数据范围内。


\end{document}


